% Second pass (Claude, 2026-11-10): first pass checked against the archive images; changes in LM_C4-039-057_changes.md.
\scanpage{39}
\mspage{667}

\[
\tag{24} l_1^\nu=\lambda_1^\nu\omega^{\nu_2}.
\]

Setting (cf. (19))
\[
\tag{25}
\left\{\begin{aligned}
\beta&=\beta_0\tau^{\nu_2}\qquad(\unsure{\beta_0|\beta}\in\Gamma\underline M)\\
\eta_0&=\beta_0\sigma(\beta_0)^{-1}
\end{aligned}\right.
\]
we shall thus have
\[
\eta\overset{\unsure{\mathrm{dfn}}}{=}
\beta\sigma(\beta)^{-1}
=\beta_0\sigma(\beta_0)\omega^{\nu_2}
=\eta_0\omega^{\nu_2}
\]
and the loop equation takes the form
of an equation in $\underline M$
\[
\tag{26}
\unsure{\zeta|\xi}=\deleted{\unsure{\beta}}\eta_0\lambda_1^\nu,
\qquad\text{i.e. \unsure{side} of (21)}
\]
\[
\tag{26bis}\unsure{\zeta|\xi}=
\]
where $\unsure{\zeta|\xi},\eta_0$ are two sections of $\underline M$ which
determine one another (for given $\nu$).
We have thus expressed in a
convenient way the loop equation, i.e. the
condition c)\inserted{of (17)}, it remains to express a) and b),
which are also formulated as
\[
\tag{27}
\left\{\begin{aligned}
\text{a)}\;&\unsure{\zeta|\xi}\ \text{locally of the form}\quad
\overset{\inserted{(\nu_3=0)}}{\alpha_0\rho(\alpha_0)^{-1}}
\quad\text{or}\quad
\overset{\inserted{(\nu_3=1)}}{\alpha_0\lambda_1^{-1}\rho(\alpha_0)^{-1}}\\
&\text{with }\alpha_0\in\Gamma\underline M\\
\text{b)}\;&\eta_0\ \text{locally of the form}\quad\beta_0\sigma(\beta_0)^{-1}
\end{aligned}\right.
\]
(Conditions whose formulation is independent
of $\nu,\nu_2$ \deleted{\unsure{for|$\nu_3$})}, but depends \unsure{apparently} on
$\nu_3$.) To make these conditions explicit, we
are going to write $\underline M$ additively, the relations
read
\margin{We have thus also set $\alpha=\alpha_0\tau^{\nu_3}$}
\[
\tag{28}
\left\{\begin{aligned}
\unsure{\zeta|\xi}&=\alpha_0-\rho(\alpha_0)
\quad\overset{\inserted{\text{case}}}{(\nu_3=0)}
\quad\text{or}\quad
\unsure{\zeta|\xi}=\alpha_0-\rho(\alpha_0)-\lambda_1
\quad\overset{\inserted{\text{case}}}{(\nu_3=1)}\\
\eta_0&=\beta_0-\sigma(\beta_0),
\end{aligned}\right.
\]
\margin{38}
\margin{38\\\texttt{https://grothendieck.umontpellier.fr}\\Université de Montpellier, UMR CNRS 5149, Case courrier 051, Place Eugène Bataillon\\34095 Montpellier cedex 5 - France}

\scanpage{40}

which we also express by
\[
\tag{29}
\left\{\begin{aligned}
\unsure{\zeta|\xi}&\in\operatorname{Im}(\mathrm{Id}_{\underline M}-\rho_{\underline M})
\quad(\nu_3=0)\quad\text{or}\quad
\unsure{\zeta|\xi}\in-\lambda_1+\operatorname{Im}(\mathrm{Id}_{\underline M}-\rho_{\underline M})\\
\eta_0&\in\operatorname{Im}(\mathrm{Id}_{\underline M}-\sigma_{\underline M}).
\end{aligned}\right.
\]

It is time to move on to explicit
matrix calculations, by referring $\underline M$ to
the basis $\lambda_0,\lambda_1$. In this basis, the
\unsure{operations} of $\rho_{\underline M},\sigma_{\underline M},\tau_{\underline M}$ are given by
\[
\tag{30}
\left\{\begin{aligned}
\rho_{\underline M}&=\begin{pmatrix}0&-1\\1&-1\end{pmatrix}
&\mathrm{id}_{\underline M}-\rho_{\underline M}&=\begin{pmatrix}1&1\\-1&2\end{pmatrix}\\
\sigma_{\underline M}&=\begin{pmatrix}0&1\\1&0\end{pmatrix}
&\mathrm{id}_{\underline M}-\sigma_{\underline M}&=\begin{pmatrix}1&-1\\-1&1\end{pmatrix}\\
\tau_{\underline M}&=\begin{pmatrix}-1&0\\0&-1\end{pmatrix}.&&
\end{aligned}\right.
\]

We have
\[
\det(\mathrm{id}_{\underline M}-\rho_{\underline M})=3,
\qquad\det(\mathrm{id}_{\underline M}-\sigma_{\underline M})=0
\]
which shows that \deleted{$\mathrm{id}_{M}$} the endomorphisms
$\mathrm{id}_{\underline M}-\rho_{\underline M}$ and $\mathrm{id}_{\underline M}-\sigma_{\underline M}$ of $\underline M$ are
not quite automorphisms\inserted{i.e. epimorphisms} ---
It screws up in characteristic 3 for the first, \uline{everywhere}
\deleted{\unsure{characteristic 2}}\inserted{(in every characteristic)} for the second.
\deleted{We shall therefore take}\inserted{We shall have}
\deleted{more} reason to take as ``parameters''
not $\unsure{\zeta|\xi}$ and $\eta_0$, but quite obviously
$\alpha_0$ and $\beta_0$ (in addition to the discrete ``parameters''
$\nu_2,\nu_3$ and $\nu'$), and to \deleted{define}\inserted{make} \deleted{thus} a
group (if possible\ldots) with the systems

\margin{39}
\margin{39\\\texttt{https://grothendieck.umontpellier.fr}\\Université de Montpellier, UMR CNRS 5149, Case courrier 051, Place Eugène Bataillon\\34095 Montpellier cedex 5 - France}

\scanpage{41}
\mspage{668}

\deleted{\unsure{(24)}}
\[
\tag{31}
\widetilde{\underline Z}_{\rho,\sigma}
=\left\{(\nu,\nu',\nu_2,\nu_3,\alpha_0,\beta_0)\ \left|\
\begin{gathered}
\nu,\nu'\in\Gamma\mathbf G_a,\quad
\nu_2,\nu_3\in\Gamma\{0,1\}_{S_0},\quad
\alpha_0,\beta_0\in\Gamma\underline M\\
\left\{\begin{aligned}
\unsure{\nu}&=2\nu'\unsure{-|+}\nu_2\qquad 2\nu+1\in\Gamma\mathbf G_m\\
\underbrace{\alpha_0-\rho(\alpha_0)\unsure{-|+}\nu_3\lambda_1}_{\unsure{\zeta|\xi}}
&=\underbrace{\beta_0-\sigma(\beta_0)}_{\eta}+\nu\lambda_1
\end{aligned}\right.
\end{gathered}\right.\right\}
\]
\margin{\uline{N.B.} $\nu'$ determines $\nu$ which is redundant}

The loop equation, in $\alpha_0,\beta_0$, thus takes
the form
\[
\tag{32}
(\unsure{\mathrm{id}|1}-\rho_{\underline M})\alpha_0
\unsure{-|=}(\mathrm{id}-\sigma_{\underline M})(\beta_0)
=\bigl(\underbrace{2\nu'\unsure{-|+}\nu_2}_{\nu}\unsure{+|-}\nu_3\bigr)\lambda_1
\]

\margin{This must indeed \unsure{come from the} fact that the homomorphism
\diagram{\begin{tikzcd}[column sep=large]
\underline M\times\underline M\arrow[r]&\underline M\\
(\alpha_0,\beta_0)\arrow[r,mapsto]&
(\mathrm{id}-\rho_{\underline M})\alpha_0-
(\deleted{\unsure{\mathrm{id}}}\inserted{1}-\sigma_{\underline M})(\beta_0)
\end{tikzcd}}
\unsure{is|is} indeed an epimorphism\ldots\ But this is the case only in characteristic $\ne3$ !}

\uline{Remark} The scheme $\underline Z_{\rho,\sigma}$ itself, as it
was defined in \S\ 49 \deleted{\unsure{\S\ 6}} $\overline{\underline{\mathrm{VI}}}$, is not smooth
it seems to me --- \unsure{it gets stuck|gets complicated} in characteristics 2 and 3.
On the other hand the finer group $\widetilde{\underline Z}_{\rho,\sigma}$,
defined by bringing in the
$\alpha,\beta$ right away (instead of $(\unsure{\zeta|\xi},\eta)$)
does seem \uline{smooth}. \unsure{This} leads once again
to modifying a little the presentation of the
constructions of \S\ 49, by \unsure{denying myself|not using}
\unsure{yet another} group $\underline Z_{\rho,\sigma}$ (tied to the coordinates
$\unsure{\zeta|\xi},\eta$), and by replacing it with
a group which \unsure{will appear|appears} \inserted{\unsure{at the same time|in a way}} as a
normalizer of $L_0$ (or of $L_0^{\unsure{\sharp}}$) in $\mathcal M_{\rho,\sigma}$
(or $\mathcal M_{\rho,\sigma}^{\unsure{\natural|\sharp}}$). We shall give a presentation along
these lines in the \deleted{\unsure{sequel}} \S\ that follows, and we
\unsure{assume} here that this construction is already in hand,



\margin{40}
\margin{40\\\texttt{https://grothendieck.umontpellier.fr}\\Université de Montpellier, UMR CNRS 5149, Case courrier 051, Place Eugène Bataillon\\34095 Montpellier cedex 5 - France}

\scanpage{42}

Whence the \deleted{\unsure{sphere|scheme}} (scheme \deleted{and) group}
$\widetilde{\underline Z}_{\rho,\sigma}$ will thus have
a suitable group structure, and
operates on $\unsure{\underline C^{+}|\underline G^{+}}$ via
\[
\tag{33}
\left\{\begin{aligned}
u(\rho)&=(\alpha_0-\rho(\alpha_0)+\nu_3\lambda_1)\rho\\
u(\sigma)&=(\beta_0-\sigma(\beta_0))\omega^{\nu_2}\sigma.
\end{aligned}\right.
\]

But we must all the same check that
these formulas (33) really do
determine an automorphism $u$ of $\unsure{\underline C^{+}|\underline G^{+}}$. The
essential check : make sure
that the \deleted{formula \unsure{[7]}}\inserted{relation} (7)
\[
\tag{34}[\sigma,\lambda_\infty]
\]
is respected by $u$.

I am going to make explicit the action of $u$ on
$\underline M$, by making explicit its values on
\[
\left\{\begin{aligned}
\lambda_0&=\omega l_0=\omega\sigma\rho\sigma\rho\\
\lambda_1&=\omega l_1=\omega\rho\sigma\rho\sigma
\end{aligned}\right.
\]
in terms of \deleted{$\unsure{\zeta|\xi}$ and $\eta_0$, \unsure{intro}}
\[
\tag{35}
\left\{\begin{aligned}
\unsure{\zeta|\xi}&=\alpha_0-\rho(\alpha_0)+\nu_3\lambda_1\\
\eta_0&=\beta_0-\sigma(\beta_0)
\end{aligned}\right.
\qquad\unsure{\alpha|\&}\qquad\in\Gamma\underline M
\]
Linked to $u$ by
\[
\tag{36}
u\rho=\unsure{\zeta|\xi}\rho,\qquad
u\sigma=\omega^{\nu_2}\eta_0\sigma.
\]
We shall have \deleted{thus} by the loop relation
\[
u(\varepsilon_0)=\deleted{\unsure{\mu}}\varepsilon_0^{\deleted{\unsure{\mu}}\inserted{(\mu)}}
\qquad(\mu=2\nu+1=4\nu'+2\nu_2+1)
\]
whence

\margin{41}
\margin{41\\\texttt{https://grothendieck.umontpellier.fr}\\Université de Montpellier, UMR CNRS 5149, Case courrier 051, Place Eugène Bataillon\\34095 Montpellier cedex 5 - France}

\scanpage{43}
\mspage{669}
\[
u(l_0)=u(\varepsilon_0^2)=l_0^{(\mu)}
\]
and
\[
u(\lambda_0)=\unsure{u}(\omega l_0)
=\underbrace{u(\omega)}_{=\omega}\,u(l_0)
=\omega l_0^{(\mu)}
=\underbrace{\omega l_0}_{\lambda_0}
 \bigl(\underbrace{l_0^{(2\nu)}}_{\lambda_0^{2\nu}}\bigr)
=\lambda_0^{(2\nu+1)}=\lambda_0^{(\mu)}
\]
as it should. (To reassure myself, I have checked that the second formula (33) does indeed give
\[
u(\sigma^2)=(\eta_0\sigma)(\eta_0\sigma)
=\bigl(\underbrace{\eta_0\,\sigma(\eta_0)}_{1}\bigr)
 \deleted{\unsure{\sigma^2}}\omega=\omega
\quad\text{i.e.}\quad\deleted{\text{and}}
\]
\[
\tag{37} u(\omega)=\omega
\]
We shall have
\[
u(\lambda_1)=u(\rho(\lambda_0))
=(\unsure{\zeta|\xi}\rho u)(\lambda_0)
=(\unsure{\zeta|\xi}\rho)u(\lambda_0)
=(\unsure{\zeta|\xi}\rho)(\lambda_0^\mu)=\lambda_1^\mu
\]
thus finally, as predicted
\[
\tag{38} u(\lambda)=\mu.\lambda\qquad\text{for}\quad\lambda\in\Gamma\underline M
\]
\margin{where $\mu=2\nu+1=4\nu'+2\nu_2+1$}
We shall indeed have that $u(\sigma)=(\beta_0-\sigma(\beta_0))\omega^{\nu_2}.\sigma$
commutes with $u(\lambda_\infty)=\mu.\lambda_\infty$.

All these calculations are a bit \unsure{murky}, logically, for at the same time I want to prove
the existence of a \deleted{\unsure{homomorphism}} $u$ \inserted{of $\unsure{\underline C^{+}|\underline G^{+}}$},
and in the preceding calculations I already assume it (\unsure{implicitly}, in a
``bigger'' group, namely the \unsure{smooth|free} group scheme \inserted{(\unsure{heuristic})}
generated by $\rho,\sigma\ldots$.
One way of making them \uline{rigorous} would be to use the intrinsic description of
$\unsure{\underline C^{+}|\underline G^{+}}$ :
\deleted{to \unsure{prove}}
\[
\tag{39}
\unsure{\underline C^{+}|\underline G^{+}}
=\unsure{C^{+}_{S_0}|G^{+}_{S_0}}
 \mathop{\wedge}_{M_{S_0}}\underline M_{S_0}
\]
thus defining a homomorphism\inserted{$u$} from
$\unsure{\underline C^{+}|\underline G^{+}}$ to an $S_0$-group $G$
(here $G=\unsure{\underline C^{+}|\underline G^{+}}$) amounts to
\margin{42}
\margin{42\\\texttt{https://grothendieck.umontpellier.fr}\\Université de Montpellier, UMR CNRS 5149, Case courrier 051, Place Eugène Bataillon\\34095 Montpellier cedex 5 - France}